\documentclass[10pt,a4paper]{article}

\usepackage[utf8]{inputenc}
\usepackage[T1]{fontenc}
\usepackage[margin=1.5cm]{geometry}
\usepackage{amsmath,amssymb}
\usepackage{xcolor}
\usepackage{tcolorbox}
\tcbuselibrary{breakable}
\usepackage{listings}
\usepackage{hyperref}
\usepackage{tikz}
\usetikzlibrary{positioning,calc,arrows.meta,shapes.geometric}

\lstset{
    language=C,
    basicstyle=\ttfamily\scriptsize,
    backgroundcolor=\color{gray!10},
    keywordstyle=\color{blue},
    commentstyle=\color{gray},
    stringstyle=\color{red},
    numbers=left,
    numberstyle=\tiny,
    breaklines=true,
    frame=single,
    showstringspaces=false,
    tabsize=4,
    inputencoding=utf8,
    extendedchars=true,
    literate={é}{{\'e}}1 {è}{{\`e}}1 {à}{{\`a}}1 {ç}{{\c{c}}}1 
             {ê}{{\^e}}1 {î}{{\^i}}1 {ô}{{\^o}}1 {û}{{\^u}}1 {ù}{{\`u}}1
             {É}{{\'E}}1 {È}{{\`E}}1 {À}{{\`A}}1 {Ç}{{\c{C}}}1
             {Ê}{{\^E}}1 {Î}{{\^I}}1 {Ô}{{\^O}}1 {Û}{{\^U}}1 {Ù}{{\`U}}1
             {ë}{{\"e}}1 {ï}{{\"\i}}1 {ö}{{\"o}}1 {ü}{{\"u}}1
             {Ë}{{\"E}}1 {Ï}{{\"I}}1 {Ö}{{\"O}}1 {Ü}{{\"U}}1
             {œ}{{\oe}}1 {Œ}{{\OE}}1 {æ}{{\ae}}1 {Æ}{{\AE}}1
}

\newtcolorbox{infobox}[1]{
    colback=blue!5,colframe=blue!75!black,
    fonttitle=\bfseries,title=#1,breakable
}

\newtcolorbox{exercicebox}[1]{
    colback=green!5,colframe=green!50!black,
    fonttitle=\bfseries,title=#1,breakable
}

\newtcolorbox{warningbox}[1]{
    colback=red!5,colframe=red!75!black,
    fonttitle=\bfseries,title=#1,breakable
}

\newtcolorbox{techbox}[1]{
    colback=orange!5,colframe=orange!75!black,
    fonttitle=\bfseries,title=#1,breakable
}

\title{\textbf{TP Chapitre 0 : Compilation et Débogage}\\
\large Du Code Source à l'Exécutable}
\author{Dr. El Hadji Bassirou TOURÉ\\
Département de Mathématiques et Informatique\\
Faculté des Sciences et Techniques\\
Université Cheikh Anta Diop de Dakar}
\date{}

\begin{document}

\maketitle

\begin{infobox}{Environnement de travail}
Ce TP doit être réalisé dans l'environnement Docker du cours. 
Si vous n'avez pas encore configuré l'environnement, consultez le TP-1 (Setup Environnement).

\textbf{Commandes de base:}
\begin{verbatim}
# Linux/macOS
docker run -it -v $(pwd):/workspace elbachirtoure/c-advanced-course:v1.0

# Windows (Git Bash)
docker run -it -v /${PWD}:/workspace elbachirtoure/c-advanced-course:v1.0

# Dans le container
cd /workspace
\end{verbatim}
\end{infobox}

\section{Phases de Compilation et Optimisations}

\begin{exercicebox}{Exercice 1 : De la source à l'exécutable}
\textbf{Objectifs :} Observer les fichiers intermédiaires, comprendre le rôle de chaque phase, mesurer l'impact des optimisations.

\textbf{Partie A -- Fichiers intermédiaires.} Créer \texttt{phases.c} :
\begin{lstlisting}
#include <stdio.h>

#define TAILLE 100
#define CARRE(x) ((x) * (x))

int global = 42;  // Variable globale initialisée

int somme_carres(int n) {
    int total = 0;
    for (int i = 1; i <= n; i++) {
        total += CARRE(i);
    }
    return total;
}

int main() {
    int local = TAILLE;
    printf("Somme des carrés de 1 à %d = %d\n", local, somme_carres(local));
    return 0;
}
\end{lstlisting}

Générer tous les fichiers intermédiaires :
\begin{lstlisting}[language=bash]
gcc -E phases.c -o phases.i      # Préprocesseur
gcc -S phases.c -o phases.s      # Compilation (assembleur)
gcc -c phases.c -o phases.o      # Assemblage (objet)
gcc phases.o -o phases           # Édition des liens
\end{lstlisting}

Analyser les fichiers :
\begin{lstlisting}[language=bash]
head -50 phases.i && tail -20 phases.i   # Fichier préprocessé
cat phases.s                              # Code assembleur
nm phases.o                               # Symboles avant linkage
nm phases                                 # Symboles après linkage
\end{lstlisting}

\textbf{Partie B -- Impact des optimisations.} Compiler avec différents niveaux :
\begin{lstlisting}[language=bash]
gcc -O0 -S phases.c -o phases_O0.s && gcc -O0 -o phases_O0 phases.c
gcc -O2 -S phases.c -o phases_O2.s && gcc -O2 -o phases_O2 phases.c
gcc -O3 -S phases.c -o phases_O3.s && gcc -O3 -o phases_O3 phases.c

wc -l phases_O*.s                # Comparer le nombre de lignes
ls -la phases_O*                 # Comparer les tailles
diff phases_O0.s phases_O2.s | head -30
\end{lstlisting}

\textbf{Questions :}
\begin{enumerate}
\item Les macros \texttt{TAILLE} et \texttt{CARRE} apparaissent-elles dans \texttt{phases.i} ? Pourquoi ?
\item Dans \texttt{nm phases.o}, que signifient les symboles marqués \texttt{T}, \texttt{D}, et \texttt{U} ?
\item La fonction \texttt{somme\_carres} est-elle présente dans l'assembleur \texttt{-O3} ? Expliquer.
\item Pourquoi le nombre de lignes assembleur peut-il augmenter avec l'optimisation ?
\end{enumerate}
\end{exercicebox}

\section{Compilation Séparée et Erreurs}

\begin{exercicebox}{Exercice 2 : Multi-fichiers et diagnostic d'erreurs}
\textbf{Objectifs :} Maîtriser la compilation séparée, distinguer erreurs de compilation et de liaison.

\textbf{Partie A -- Compilation séparée.} Créer trois fichiers :

\texttt{math\_utils.h} :
\begin{lstlisting}
#ifndef MATH_UTILS_H
#define MATH_UTILS_H
int addition(int a, int b);
int factorielle(int n);
extern int nb_appels;
#endif
\end{lstlisting}

\texttt{math\_utils.c} :
\begin{lstlisting}
#include "math_utils.h"
int nb_appels = 0;

int addition(int a, int b) { nb_appels++; return a + b; }

int factorielle(int n) {
    nb_appels++;
    return (n <= 1) ? 1 : n * factorielle(n - 1);
}
\end{lstlisting}

\texttt{main.c} :
\begin{lstlisting}
#include <stdio.h>
#include "math_utils.h"

int main() {
    printf("5 + 3 = %d\n", addition(5, 3));
    printf("5! = %d\n", factorielle(5));
    printf("Nombre d'appels: %d\n", nb_appels);
    return 0;
}
\end{lstlisting}

Compiler et analyser :
\begin{lstlisting}[language=bash]
gcc -Wall -c math_utils.c -o math_utils.o
gcc -Wall -c main.c -o main.o
nm math_utils.o && echo "---" && nm main.o
gcc math_utils.o main.o -o programme
./programme
\end{lstlisting}

\textbf{Partie B -- Provoquer des erreurs.} Tester ces scénarios :
\begin{lstlisting}[language=bash]
# Erreur de compilation : oublier un point-virgule dans main.c
# Erreur de linkage : compiler sans math_utils.o
gcc main.o -o test_link

# Définition multiple : ajouter "int nb_appels = 99;" dans main.c
gcc -c main.c -o main.o   # Compile OK
gcc math_utils.o main.o -o test_multi  # Erreur ici
\end{lstlisting}

\textbf{Questions :}
\begin{enumerate}
\item Pourquoi \texttt{main.o} montre \texttt{addition} comme \texttt{U} (undefined) alors que le code compile ?
\item Quelle est la différence entre le message d'erreur du compilateur et celui du linker ?
\item Comment le linker détecte-t-il une définition multiple de \texttt{nb\_appels} ?
\end{enumerate}
\end{exercicebox}

\section{Débogage avec GDB}

\begin{exercicebox}{Exercice 3 : Trouver et corriger des bugs avec GDB}
\textbf{Objectifs :} Utiliser breakpoints, inspection de variables, exécution pas à pas, pile d'appels.

Créer \texttt{buggy.c} contenant plusieurs bugs :
\begin{lstlisting}
#include <stdio.h>

int chercher_max(int tab[], int taille) {
    int max = 0;  // Bug 1: mauvaise initialisation
    for (int i = 0; i <= taille; i++) {  // Bug 2: <= au lieu de <
        if (tab[i] > max)
            max = tab[i];
    }
    return max;
}

double moyenne(int tab[], int taille) {
    int somme = 0;
    for (int i = 0; i < taille; i++)
        somme += tab[i];
    return somme / taille;  // Bug 3: division entière
}

int main() {
    int valeurs[] = {-5, -2, -8, -1, -3};
    int n = 5;
    
    printf("Maximum: %d\n", chercher_max(valeurs, n));  // Attendu: -1
    
    int notes[] = {15, 12, 18, 10, 14};
    printf("Moyenne: %.2f\n", moyenne(notes, 5));  // Attendu: 13.80
    
    return 0;
}
\end{lstlisting}

Compiler et déboguer :
\begin{lstlisting}[language=bash]
gcc -g -O0 -Wall buggy.c -o buggy
./buggy  # Observer les résultats incorrects

gdb ./buggy
(gdb) break chercher_max
(gdb) run
(gdb) print taille
(gdb) print tab[0]
(gdb) watch max
(gdb) continue       # Observer les changements de max
(gdb) print i        # Quand i == taille, accès hors limites
(gdb) backtrace
(gdb) break moyenne
(gdb) continue
(gdb) print somme
(gdb) print somme / taille    # Division entière!
(gdb) print (double)somme / taille
(gdb) quit
\end{lstlisting}

\textbf{Questions :}
\begin{enumerate}
\item Pourquoi \texttt{chercher\_max} retourne 0 pour un tableau de valeurs négatives ?
\item Quelle valeur lit \texttt{tab[5]} (accès hors limites) ? Utilisez GDB pour le vérifier.
\item Comment corriger la division entière dans \texttt{moyenne} ?
\item Corrigez les trois bugs et vérifiez les résultats attendus.
\end{enumerate}
\end{exercicebox}



\section*{Livrables}

\begin{infobox}{À rendre}
\begin{enumerate}
\item Un fichier pdf en bonne et due forme contenant les captures, réponses et explications
\end{enumerate}

\end{infobox}

\end{document}
